%! ~~~ Packages Setup ~~~ 
\documentclass[]{../../../../tex/classes/styledArticle}
\usepackage{../../../../tex/packages/hebrewSupport}
\usepackage{../../../../tex/packages/mathShortcuts}
\usepackage{../../../../tex/packages/theoremsSupport}
\usepackage{../../../../tex/packages/enfancysections}

%! ~~~ Document ~~~

\author{שחר פרץ}
\title{\textit{לינארית 2א 14}}
\date{26 במאי 2025}
\begin{document}
    \maketitle
    
    \defi{יהיו $A, A' \in M_n(\F)$, נאמר שהן חופפות אם קיימת הפיכה $P \in M_n(\F)$ כך ש־$A' = P^TAP$. }
    \theo{מטריצות חופפות אמ''מ הן מייצגות את אותה התבנית הבילינארית. }
    \theo{אם $A, A'$ חופפות, אז: 
    \begin{enumerate}
        \item \hfil $\rk A = \rk A^T$
        \item \hfil $\exists 0 \neq c \in \F \co \det A' = c^2 \det A$
    \end{enumerate}}
    \begin{proof}
        הגדרנו $\rk f$ כאשר $f$ בילינ' להיות הדרגה של המייצגת את התבנית, וראינו שהיא לא תלויה בבסיס. בכך למעשה כבר הוכחנו את \textit{1.}. עבור \textit{2.}, מתקיים $A' = P^TAP$ ו־$P$ הפיכה (ולמעשה מטריצת מעבר בסיס) ולכן אם נסמן $c = |P| = |P^T|$ מתקיים: 
        \[ |A'| = |P^TAP| = |P^T| \cdot |A| \cdot |P| = c^2 |A| \]
    \end{proof}
    (הערה: יש שדות שמעליהם טענה 2 לא מעניינת במיוחד). 
    
    \defi{תבנית $f$ מעל $V$ נקראת \textit{סימטרית} אם \hfill $\forall v, w \in V \co f(v, w) = f(w, v)$}
    \defi{תבנית $f$ מעל $V$ נקראת \textit{אנטי־סימטרית} אם \hfill $\forall v, w \in V \co f(v, w) = -f(w, v)$}
    נבחין שאם $\chr \F \neq 2$, ניתן להגדיר את: 
    \[ \phi(v, w) = \frac{f(v, w) + f(w, v)}{2}, \ \psi(v, w) = \frac{f(v, w) - f(w, v)}{2} \]
    מתקיים ש־$\phi$ סימ' ו־$\psi$ א־סימ' וכן $f = \phi + \psi$. 
    
    \theo{תהי $f$ תבנית בילינ' על $V$, ו־$B = (v_i)_{i = 1}^{n}$ בסיס ל־$B$. נניח $A =(a_{ij})^{n}_{i, j = 1}$ המייצגת את $f$ ביחס ל־$B$. אז $f$ סימ'/אסימ' אממ $A$ סימ'/אסימ'. }
    
    \begin{proof}
        \begin{itemize}
            \item[$\implies$] אם $f$ סימ'/אסימ', אז: 
            \begin{align*}
                a_{ij} = f(v_i, v_j) &\overset{\Sym}{=} f(v_j, v_i) = a_{ji} \\
                a_{ij} = f(v_i, v_j) \,\,&\!\!\overset{\Asym}{=} -f(v_j, v_i) = -a_{ji}
            \end{align*}
            \item[$\impliedby$] אם $A$ סימ' אז: 
            \[ f(v, w) = [u]_B^TA[w]_B \overset{(1)}{=} ([u]_B^TA[w]_B)^T = [w]_B^TA^T([u]_B^T)^T = [w]^T_BA[u]_B = f(w, v) \]
            כאשר $(1)$ מתקיים כי transpose למטריצה מגודל $1 \times 1$ מחזיר אותו הדבר. וכן במקרה האנטי־סימטרי: 
            \[ f(u, w) = [w]^T_B(-A)[u]_B = -[w]_B^TA[u]_B = -(w, u) \]
        \end{itemize}
    \end{proof}
    
    תת־פרק חדש: 
    \section{\en{Quadratic form}}
    \defi{תהא $f$ תבנית על $V$. התבנית הריבועית: 
    \[ Q_p \co V \to \F, \ Q_f(v) = f(v, v) \]}
    היא \textit{לא} העתקה לינארית. היא תבנית ריבועית. \textbf{דוגמאות: }
    \begin{itemize}
        \item 
        \[ f\cl{\pms{x \\ y}, \ \pms{u \\ v}} = xv + uy \implies Q\cl{\pms{x \\ y}} = xy + yx = 2xy \]
        \item 
        \[ f\cl{\pms{x \\ y}, \ \pms{u \\ v}} = xv - uy \implies Q\cl{\pms{x \\ y}} = xy - yx = 0 \]
        \item התבנית הסטנדרטית: 
        \[ f\cl{\pms{x \\ y}, \pms{u \\ v}} = xu + yv \implies Q_f\pms{x \\ y} = x^2 + y^2 \]
    \end{itemize}
    
    \noti{עבור תבנית בילינארית $f$ על $V$, נגדיר את $\h f(u, v) = f(v, u)$}
    אם $f$ סימטרית נבחין ש־$Q_f = Q_{\h f}$
    
    \theo{תהי $f$ תבנית בילי' סימ' על $V$, ונניח ש־$\chr \F \neq 2$, אז: 
    \begin{itemize}
        \item \hfil $f(v, w) = \frac{Q_f(v + w) - Q_f(v) - Q_f(w)}{2}$
        \item אם $f$ איינה תבנית ה־$0$ אז קיים $0 \neq v \in V$ כך ש־$Q_f(v) \neq 0$. 
    \end{itemize}}
    \begin{proof}
        \begin{align*}
            Q_f(v + w) - Q_f(v) - Q_f(w) = &f(v + w, v + w) - f(v, v) - f(w, w) \\
            = &\quad\, f(v, v) + f(v, w) \\
            &-f(w, v) + f(w, w) \\
            &-f(v, v) - f(w, w) \\
            \overset{\Sym}{=} \!\!&\,\, 2f(v, w)
        \end{align*}
    \end{proof}
    עבור 1, עתה נוכיח את 2: נניח $\forall v \in V \co Q_f(v) = 0$ אז
    \[ \forall v, w, \in V \co f(v, w) = \frac{Q_f(v + w) - Q_f(v) - Q_f(w)}{2} = 0  \]
    
    למה שונה ממציין 2 חשוב: 
    \[ \textstyle f(\binom{x}{y}, \ \binom{u}{v}) = xv + yu \implies Q_f = 0 \land f \neq 0 \]
    
    '
    \textit{הערה: }אין ממש טעם להגדיר תבנית ריבועית על תבנית בילינארית שאיננה סימטרית. הסיבה? כל מטריצה יכולה להיות מפורקת לחלק סימטרי וחלק אנטי־סימטרי, החלק האנטי־סימטרי יתאפס (אלכסון אפס) ורק החלק הסימטרי ישאר בכל מקרה. זה לא שאי־אפשר, זה פשוט לא מעניין במיוחד. 
    
    \theo{נניח $\chr F \neq 2$, ו־$f$ סימטרית על $V$. אז קיים בסיס ל־$V$ הוא $B = (v_i)_{i = 1}^{n}$ כך ש־$[f]_B$ אלכסונית. }
    
    \textit{תזכורת: }$[f]_B$ סימון המוגדר בסיכום זה בלבד. בקורס מדברים על המטריצה המייצגת של בי־לינארית במילים. 
    
    \begin{proof}
        באינדוקציה על $n$. בסיס $n = 1$ ברור. אם $f$ תבנית ה־$0$, אז כל בסיס שנבחר מתאים. אחרת, קיים $0 \neq v \in V$ כך ש־$Q_f(v) \neq 0$. נגדיר $U = \{u \in V \mid f(u, v) = 0\}$. תמ''ו כי גרעין של ה''ל (כי קיבענו את $v$). מה התמונה של ההעתקה? $f(v, v) = Q_f(v) \neq 0$. לכן תמונת ההעתקה היא כל $\F$, וממדה $1$. ידוע $U$ תמ''ו מממד $n - 1$. אז $f_{|U} \co U\times U \to \F$ לכסינה ולכן קיים בסיס $B_U$ כך ש־$[f_{|U}]$ אלכסונית. נגדיר את $B = \{v\} \cup B_U$ נבחין שהיא מהצורה: 
        \[ \pms{* & 0& \cdots \\ 0 & [f_{|U}]_B& - \\ \vdots & \vert} \]
    \end{proof}
    
    \ndoc
    
\end{document}
